\section{Introduction}
% \com{Note to self: will need to highlight our unique contribution compared to other studies that have relied on the same data}

Over recent years, identity theft (IDT) has evolved from a localized crime of physical opportunity to a highly organized and technology-driven one. Where offenders once relied on stolen mail or discarded documents, the modern threat environment is characterized by the digitization of financial and medical records, which has created more opportunities for exploitation. Today, large-scale data breaches occur with such frequency that they have become a predictable feature of the digital economy \cite{Verizon_DBIR_2025}, feeding a dark-web marketplace where personal identifiers are traded as low-cost commodities \cite{Holt2013}. This availability of data means that even a small corporate data compromise can expose millions of individuals' data. Despite this growing prevalence, current reports on the economic impact of data breaches remain skewed by their focus on the corporation. When a breach occurs, industry reports typically quantify the impact in terms of corporate expenses such as regulatory fines, legal settlements, and forensic investigations \cite{IBM_DBR_2025}. However, these figures completely ignore the externalized (social) costs shifted onto the victims. 

This study stems from the observation that while a company may report a fixed dollar amount per compromised record, that figure does not account for the social harm that victims face. As criminal tactics become more and more sophisticated, the burden on a victim has shifted from temporary financial inconvenience to a long struggle for digital recovery. In some cases, victims suffer serious physical and mental ailments as a result of the crime \cite{Golladay2017Consequences}.
%
Admittedly, there are other forms of externalized cost of data breaches beyond costs associated with financial crimes suffered by consumers: for instance, data breaches may lead a firm to increase spending in cybersecurity, which raises the cost of products and services it offers, which is then at least partially transferred to its customers. In the present study we will limit ourselves to the cost of IDTs, and our goal is to develop a method that can help us estimate the totality of this cost incurred by a major data breach event. 

Toward that end, we will attempt to broaden the typical definition of the cost of an IDT, which is often narrowly measured by a victim's immediate out-of-pocket (OOP) loss, a measure that ignores or underestimates other costs that victims often incur. Quantifying a broader array of direct as well as indirect costs associated with IDT, which includes the hours victims spend on the phone with banks and government agencies, the professional legal help they may need to hire, and the physical and emotional distress that manifests as medical bills, is critical for several reasons. First, a comprehensive cost model allows policymakers to accurately weigh the benefits of security mandates against the burden taken on by the public. Second, understanding these costs helps victim service organizations and insurance providers better allocate resources for the intangible harms, such as mental health support, that traditional financial compensation ignores. By focusing on the victim's perspective, our work provides a more complete insight into the total burden of these breach events. 
\rev{It should be noted that while we substantially broaden the array of downstream costs associated with an IDT in our study, there are clearly additional externalized and downstream harm to both individuals and societies that our model does not capture. For this reason, while we will continue to use the term ``social cost'' throughout this paper for lack of a better term, we acknowledge that this term is used in a narrow sense constrained by what we can extract from available data.}



% Furthermore, current reports on the cost of a data breach are skewed by their focus on corporate liability. When a breach occurs, industry reports typically quantify the impact in terms of corporate expenses such as regulatory fines, legal settlements, and forensic investigations \cite{IBM_DBR_2025}. However, these figures ignore the externalized (social) costs shifted onto the victims themselves. While a company may report a fixed dollar amount per compromised record, that figure does not account for the social harm that victims face. 


%\com{This is a very nice paragraph, and connects the cost of IDT to the social cost of data breaches.  So perhaps we don't need to do anything else but can simply strike the word ``social'' whenever we talk about the cost of IDT, but keep this word whenever we talk about data breaches...}

%\resp{Okay, that sounds good. I took 'Social' out of the title, and we can keep using it in the paper when we refer to data breaches.}
%\resp{I am actually not sure about the title now, I think we should change it. I'm open to suggestions for a new one if you have any, otherwise I will think more about it.}

Establishing this ``social cost'' requires linking two disconnected types of data: corporate data breaches and the private experiences of IDT victims. In this paper, we bridge this gap by conducting a longitudinal analysis on trends in IDT. We develop a model to quantify the social cost of IDT, and analyze the number of IDT victims compared to the number of data records breached over time. We use over a decade of national survey data to map the life cycle of IDT, from how it is stolen, how it is discovered, and most importantly, what it truly costs the American public in both time and well-being. By pairing this victim data with the chronology of major data breaches, we test the question: \textbf{Does the frequency of mega-breaches drive a measurable increase in IDT, and if so, does the corporate penalty match the social cost?} To address this, we propose two estimation methods: an empirical lower bound of a mega-breach's short-term social cost (a lower bound), and a time-decaying projection of its long-term impact (an upper bound).
Our contributions are as follows:

\begin{enumerate}
    \item \textbf{Longitudinal data integration}: To our knowledge, this is the first study to harmonize and pool all six waves of the ITS (2008-2021) to provide a comprehensive 13-year analysis of identity theft trends.
    \item \textbf{Social cost modeling:} Unlike previous studies that focus primarily on direct financial loss, our model incorporates social costs, including the opportunity cost of lost time and monetized healthcare expenses related to mental and physical distress.
    \item \textbf{Data breach to identity theft victim analysis:} By comparing self-reported survey data with external breach chronology data, we conduct a Wilcoxon signed-rank test to verify whether there is an increase in identity theft cases following a mega-breach event, and calculate a breach to identity theft victim conversion rate to quantify the risk of successful victimization following a breach event. 
    \item \textbf{Corporate liability case studies:} We apply our social cost model and breach-to-victim conversion in three case studies: the 2009 Heartland Payment Systems breach, the 2013 Target breach, and the 2017 Equifax breach. We calculate upper and lower bound estimates for the social cost of these events, and compare that to their settlements.
\end{enumerate}

The remainder of this paper is organized as follows. Section \ref{sec:background} describes related work and the motivation for our study. In Section \ref{sec:data_and_sample}, we describe our data sources and pre-processing procedures.  Section \ref{sec:comprehensive_cost_of_IDT} details our social cost model along with our calculations and results from the ITS dataset. Section \ref{sec:breach_to_victim} compares the PRC data breach chronology data with the  victim data from the ITS survey while also detailing our model for estimating the breach-to-victim conversion rate. Section \ref{sec:breach_to_victim} also tests the relationship between mega-breaches and the number of victims, and calculates the social cost of our three case study breaches. Finally, Section \ref{sec:discussion} discusses limitations and extensions of this study, and \ref{sec:conclusion} concludes the paper. Further technical details regarding the specific data cleaning protocols, social cost calculations, and extended results are provided in the Appendix. 

% \com{Maybe we can add a line here stating many of the details on the datasets are included in the appendix or else can be found in the extended version, etc.} \resp{Done.}

% Section \ref{sec:background} describes related work and the motivation for our work. In Section \ref{sec:data_and_sample}, we describe our data sources and cleaning procedures, and in Section \ref{sec:methods}, we describe our social cost model as well as breach to identity theft victim conversion rate calculation. Our results are organized into three parts. The first part is Section \ref{sec:demographics}, which details an overview of our data and sample used for analysis, as well as the data trends for identity theft. The second part is Section \ref{sec:social_cost_identity}, which details the social cost of identity theft, and finally, the third part is Section \ref{sec:data_breach}, which displays the results of comparing number of exposed records to the number of identity theft victims, and the link between the two. Finally, Sections \ref{sec:limitations} and \ref{sec:conclusion} include a discussion of this study's limitations, and conclusion.



\section{Background and Literature Review}
\label{sec:background}

Consistent with the Bureau of Justice Statistics (BJS), we define IDT as the attempted or successful misuse of an existing account, fraudulent opening of a new account, or the misuse of personal information for other fraudulent purposes \cite{HarrellThompson2024, Harrell2024Stats}. To study the social cost of  and trends of theft over the years, we used data from the \textit{National Crime Victimization Survey (NCVS): Identity Theft Supplement (ITS)} \cite{ICPSR26362, ICPSR34735, ICPSR36044, ICPSR36829, ICPSR37923, ICPSR38501}. The ITS is a large scale survey that collects data about individuals' personal experiences with IDT directly from a nationally representative sample of U.S. households. This survey is sponsored by the U.S. Bureau of Justice Statistics (BJS), and is administered periodically as a supplement to the main NCVS survey. The ITS gathers detailed information from IDT victims, including the types of personal information compromised, how they discovered the theft, personal and financial consequences, and the actions taken in response to the theft. In this study, we use data from six waves of the ITS, conducted in 2008, 2012, 2014, 2016, 2018, and 2021 \cite{ICPSR26362, ICPSR34735, ICPSR36044, ICPSR36829, ICPSR37923, ICPSR38501}. Complimenting this, we use data breach chronology data from Privacy Rights Clearinghouse (PRC) to obtain the number of breaches and exposed records over time, and compare with the number of victims from the ITS dataset \cite{PRC_Chronology}.

Previous research using the ITS dataset has consistently identified a target suitability profile for certain types of IDT. \citet{Nevin2025} and \citet{Copes2010} found that credit card victimization risk is generally higher among individuals who are older, white, and possess higher income and education; groups likely to possess higher credit limits and more complex financial activities. While these demographics are frequently targets of existing account fraud, research shows that lower income, younger, and minority groups are disproportionately victimized by more damaging forms of IDT such as existing bank account fraud and new account fraud. \citet{Reynolds2021Differential} and \citet{DeLiema2021} further explored these disparities, demonstrating that socio-economic status significantly impacts the likelihood and severity of out-of-pocket (OOP) losses. Beyond demographic analysis, studies such as \cite{Hu2021Forecasting} have utilized machine learning on the ITS data to forecast victimization and evaluate preventative actions. Notably, most existing studies only pool a subset of the ITS data, such as three or four waves. To the best of our knowledge, our study is the first to perform analysis of all six waves, which provides a more complete picture of trends. 

% \citet{Nevin2025} used a latent class analysis pooled over four waves of the ITS survey (20012-2018) to identify victim subgroups. They noted that higher income individuals are more likely to be victims of credit card misuse, which often results in the least severe long-term consequences \cite{Nevin2025}. Similarly, \citet{Copes2010} demonstrated that while white, higher income, and highly educated individuals are most frequently victims of existing credit card fraud, women, Black individuals, and younger, lower income groups are disproportionately victimized by more damaging forms of IDT such as existing bank account fraud and new account fraud. \citet{Reynolds2021Differential} further explored these demographic disparities and studied how these characteristics influence the likelihood of incurring an OOP loss. Using logistic regression, this study found that income, education, and marital status significantly impact the odds of suffering financial loss. For instance, Hispanic/Latinx victims showed a higher likelihood of OOP loss in certain categories compared to white victims \cite{Reynolds2021Differential}. \citet{DeLiema2021} further supported these findings, demonstrating that low socio-economic status (SES) victims are more likely to lose more, with older adults at or below 150\% of the federal poverty level being more likely to suffer OOP costs \cite{DeLiema2021}. In addition, several other strategies have been used to analyze the ITS dataset. For example,  \citet{Hu2021Forecasting} used machine learning to forecast victimization and evaluate preventative actions such as checking credit reports. These existing studies pool only a subset of the ITS data. For example, \cite{Hu2021Forecasting} used three waves, and \cite{Nevin2025} used four waves. \rev{To the best of our knowledge, our study is the first to analysis of all six waves, which provides a more complete picture of trends.} 
%and allows us to} verify if the previously reported trends regarding demographics and likelihood factors still hold. 
% \com{I feel this paragraph can perhaps be shortened a little bit since our focus is now less on the demographical differences of costs and so prior work in this area is also less relevant.} \resp{See above edits.}

Our work also expands the literature on the total cost of IDT from a victim's point of view. For instance, \citet{Miller2021} %estimates the costs of various crimes in 2017, 
%This paper 
uses the 2016 ITS survey as the source of its cost estimate of IDT, summing the victims' OOP losses and lost time costs, yet did not consider medical or mental health costs. \citet{Anderson2013Measuring} provides a measure for decomposing cybercrime impact into direct, indirect, and defense costs, but they also note that their model ``generally disregard[s] distress,'' because it is difficult to measure \cite{Anderson2013Measuring}. Similarly, \cite{Romanosky2016Examining} offers an empirical estimation of the costs associated with data breaches and security incidents from a corporate perspective, but notes that their calculations do not include intangible costs such as distress and lost time \cite{Romanosky2016Examining}. By contrast, our study considers a more comprehensive measure of the cost of IDT that includes intangible harms such as monetized health expenses.

Finally, the relationship between data breaches and individual victimization is an integral part of this cost estimate that connects corporate responsibility with (socialized) individual/consumer harm; interestingly, this appears to be a relatively under-explored area. BJS research indicates that victims of IDT are twice as likely as non-victims to have been notified that their data has been exposed in a breach in the past year \cite{Harrell2024Stats}. \citet{Bisogni2020More} use 13 years (2005-2017) of United States incident data and Bayesian modeling to demonstrate that while data breaches and IDT are correlated, their relationship is largely influenced by state population size. However, beyond these general correlations, explicit breach-to-victim calculations that quantify the probability of a leaked record resulting in a successful crime are hard to come by. This stems from the inherent difficulty of these calculations. Corporate breach notifications rarely track downstream victimization, and anonymous victimization surveys lack the data to trace fraud back to a specific exposure event. Consequently, the relationship between a corporate data breach and a consumer's loss remains obscured by the messiness of attribution.

For our purpose of establishing the link between data breaches and IDT victimization, we use PRC's data breach chronology and supplement it with data from corporate disclosures and federal breach portals. 
%\com{What is the 2nd half of the preceding sentence referring to?} \resp{It's referring to how we added in data from NH and ME to the PRC data, as well as health breach data from HHS. }
By adopting and building on the framework developed in \citet{Graves2018Should} for estimating national record exposure using linear regression to extrapolate from state-level findings when national data is missing or undisclosed, we calculate a conversion rate that measures how effectively criminals convert compromised data into successful IDT.
